% id: 147-09
% book: 베이직쎈 미적분1 · 09 정적분의 활용
% section: 유형 094 움직인 거리의 활용
% figure: none
% answer: 
% answer_source: 
% variant_level: 0 (원본 전사)
% 이 파일은 build.mjs 가 items.json 에서 만든다. 고칠 때는 items.json 을 고친다.
\smash{\raisebox{1.5mm}{\dmkichul{베이직쎈 미적분1 147쪽 09번}}}
직선 선로를 $a\mkern3mu \mathrm{m/s}$의 속도로 달리는 열차에 제동을 건 지 $t$초 후의 속도 $v(t)\mkern3mu \mathrm{m/s}$가 $v(t)=a-3t$ $(0\le t\le 10)$이다. 제동을 건 후 열차가 정지할 때까지 달린 거리가 $150\mkern3mu \mathrm{m}$일 때, 양수 $a$의 값을 구하시오.
